% BHU PYQ -- Transcribed & Corrected
% Paper: HISMD11: Introduction to Tribal History of India
% Exam: Under Graduate Multidisciplinary Course Semester I Examination 2024-25 (NEP)
% Ref Code: 24220 MBT 038
% Category: Multidisciplinary (MDC)
% Department: Department of History, Faculty of Social Sciences

\documentclass[11pt,a4paper]{article}
\usepackage[a4paper,top=2.5cm,bottom=2.5cm,left=2.8cm,right=2.8cm,headheight=14pt]{geometry}
\usepackage{amsmath,amssymb}
\usepackage[T1]{fontenc}
\usepackage[utf8]{inputenc}
\usepackage{lmodern,microtype,fancyhdr,enumitem,hyperref,lastpage,parskip}

\hypersetup{
    colorlinks=true,
    urlcolor=black,
    linkcolor=black,
    pdftitle={HISMD11: Introduction to Tribal History of India -- BHU Sem I 2024-25 (NEP MDC)}
}

\pagestyle{fancy}
\fancyhf{}
\renewcommand{\headrulewidth}{0.4pt}
\renewcommand{\footrulewidth}{0.4pt}
\fancyhead[L]{\small \textit{Banaras Hindu University}}
\fancyhead[R]{\small \textit{HISMD11: Tribal History of India (Sem I 2024-25)}}
\fancyfoot[C]{\small Page \thepage\ of \pageref{LastPage}}

\newlist{parts}{enumerate}{2}
\setlist[parts,1]{label=(\alph*),leftmargin=*,itemsep=4pt,topsep=2pt}

\newcommand{\pts}[1]{\hfill \textit{[#1]}}

\begin{document}

\begin{center}
  {\large\bfseries BANARAS HINDU UNIVERSITY}\\[2pt]
  {\normalsize Under Graduate Multidisciplinary Course}\\[2pt]
  {\normalsize Semester I Examination 2024-25 (NEP)}\\[6pt]
  \rule{0.8\linewidth}{0.5pt}\\[6pt]
  {\large\bfseries HISTORY}\\[4pt]
  {\large\bfseries Paper Code: HISMD11 : Introduction to Tribal History of India}\\[2pt]
  {\small Ref. No.: 24220 MBT 038}\\[4pt]
  {\normalsize Time: 2:30 Hours \hfill Maximum Marks: 70}
\end{center}

\rule{\linewidth}{0.4pt}

\smallskip

\begin{center}
  {\large\bfseries SECTION A / खण्ड-अ}
\end{center}

\noindent \textit{Note: Answer each question. Each question carries 2 marks.}\\[2pt]
\noindent \textit{प्रत्येक प्रश्न का उत्तर दें। प्रत्येक प्रश्न 2 अंकों का है।}

\bigskip

\noindent \textbf{Question 1.} Write a short note on the following / निम्नलिखित पर संक्षिप्त टिप्पणियाँ लिखिए:\pts{10}

\begin{parts}
  \item Kol tribe\\[1pt]
        कोल जनजाति \pts{2}

  \item Any two Articles of the Constitution related to the Tribes\\[1pt]
        जनजातियों से संबंधित संविधान के कोई दो अनुच्छेद \pts{2}

  \item Ethnography\\[1pt]
        नृवंशविज्ञान \pts{2}

  \item G.S. Ghurye\\[1pt]
        जी.एस. घुर्ये \pts{2}

  \item When and why is Tribal Pride Day celebrated?\\[1pt]
        जनजातीय गौरव दिवस कब और क्यों मनाया जाता है? \pts{2}
\end{parts}

\bigskip
\rule{\linewidth}{0.2pt}
\bigskip

\begin{center}
  {\large\bfseries SECTION B / खण्ड-ब}
\end{center}

\noindent \textit{Note: Answer each question. Each question carries 10 marks.}\\[2pt]
\noindent \textit{प्रत्येक प्रश्न का उत्तर दें। प्रत्येक प्रश्न 10 अंकों का है।}

\bigskip

\noindent \textbf{Question 2.} Write a short note on tribes of India and their general characteristics.\\[1pt]
भारत की जनजातियों और उनकी विशेषताओं पर एक संक्षिप्त टिप्पणी लिखें। \pts{10}

\begin{center}
  \textbf{OR / अथवा}
\end{center}

\noindent Elucidate on the geographical \& demographic distribution of tribes in India.\\[1pt]
भारत में जनजातियों के भौगोलिक और जनसांख्यिकीय वितरण पर प्रकाश डालें। \pts{10}

\medskip\rule{\linewidth}{0.2pt}\medskip

\noindent \textbf{Question 3.} Discuss tribes as defined by western anthropologists.\\[1pt]
पश्चिमी मानवविज्ञानियों द्वारा परिभाषित 'जनजाति' पर चर्चा करें। \pts{10}

\begin{center}
  \textbf{OR / अथवा}
\end{center}

\noindent Discuss the various sources for the study of Indian tribes.\\[1pt]
आदिवासी इतिहास के विभिन्न स्रोतों की चर्चा करें। \pts{10}

\medskip\rule{\linewidth}{0.2pt}\medskip

\noindent \textbf{Question 4.} What was the condition of tribes in colonial India?\\[1pt]
औपनिवेशिक भारत में जनजातियों की स्थिति क्या थी? \pts{10}

\begin{center}
  \textbf{OR / अथवा}
\end{center}

\noindent Discuss \textit{any one} tribal movement during the colonial period.\\[1pt]
औपनिवेशिक काल के दौरान किसी एक आदिवासी आंदोलन की चर्चा करें। \pts{10}

\bigskip
\rule{\linewidth}{0.2pt}
\bigskip

\begin{center}
  {\large\bfseries SECTION C / खण्ड-स}
\end{center}

\noindent \textit{Note: Each question carries 15 marks.}\\[2pt]
\noindent \textit{प्रत्येक प्रश्न 15 अंकों का है।}

\bigskip

\noindent \textbf{Question 5.} Write an essay on the tribes of ancient India with special focus on the Hindu epics.\\[1pt]
हिंदू महाकाव्यों पर विशेष ध्यान देते हुए प्राचीन भारत की जनजातियों पर एक निबंध लिखें। \pts{15}

\begin{center}
  \textbf{OR / अथवा}
\end{center}

\noindent Discuss the problems of the tribes arising out of industrialisation.\\[1pt]
औद्योगिकीकरण से उत्पन्न आदिवासियों की समस्याओं पर चर्चा करें। \pts{15}

\medskip\rule{\linewidth}{0.2pt}\medskip

\noindent \textbf{Question 6.} Discuss the various means of tribal economic subsistence.\\[1pt]
आदिवासी जीवन-निर्वहन के विभिन्न आर्थिक साधनों की चर्चा करें। \pts{15}

\begin{center}
  \textbf{OR / अथवा}
\end{center}

\noindent Discuss the various tribal welfare measures in post-independence India.\\[1pt]
स्वतंत्र्योत्तर भारत में आदिवासी कल्याण की विभिन्न योजनाओं की चर्चा करें। \pts{15}

\vfill
\rule{\linewidth}{0.4pt}

\begin{center}\small
\href{https://bhu-exam.vercel.app/}{Click here to visit our website}\\[4pt]
\href{https://me-aryan.vercel.app/}{Click here to visit our contributor}
\end{center}

\end{document}
